\documentclass[10pt,a4paper]{amsart}
\usepackage[margin=19mm]{geometry}
\usepackage{amsmath,amssymb,mathtools}
\usepackage[scr=rsfs,cal=euler]{mathalfa}
\usepackage{xfrac}
\usepackage[unicode,hidelinks,bookmarks=false]{hyperref}
\newcommand{\ie}{i.e.\ }
\newcommand{\cf}{cf.\ }
\newcommand{\resp}{respectively\ }
\newcommand{\Subsection}{\subsection}
\providecommand{\nobreakdash}{\nobreak\hbox{-}}
\setlength{\emergencystretch}{2em}
\multlinegap0pt
\usepackage{mathtools}
\usepackage{amssymb}
\newtheorem{lemma}{Lemma}
\newtheorem{proposition}{Proposition}
\newtheorem{claim}{Claim}
\def\L{\mathcal{L}}
\def\ZZ{\mathbb{Z}}
\def\g{\mathfrak{g}}
\def\kk{\Bbbk}
\newcommand\restr[2]{{\left.\kern-\nulldelimiterspace #1 \vphantom{\big|} \right|_{#2}}}
\DeclareMathOperator{\spn}{span}
\title{Lie subalgebras of vector fields on curves}
\author{Lucas Buzaglo, Colin Ingalls}
\date{Source: arXiv:2606.21748v1 (CC BY 4.0)}
\begin{document}
\maketitle

\noindent\emph{Source and licence.} Lie subalgebras of vector fields on curves. Version arXiv:2606.21748v1 (CC BY 4.0), distributed by arXiv under the Creative Commons Attribution 4.0 International licence, http://creativecommons.org/licenses/by/4.0/, as recorded in the arXiv licence metadata for this exact version and shown on the abstract page https://arxiv.org/abs/2606.21748v1. Full licence notice: \texttt{licenses/arxiv-2606.21748-CC-BY-4.0.txt}.\par\smallskip

\noindent\textit{Editorial scope.} Counted unit: the article's proposition q (source0.tex lines 699--705). The article's complete original proof of it is delivered with it (source0.tex lines 706--809, one \texttt{proof} environment of 104 lines). No mathematics is written, repaired or re-derived: the statement and the proof are the article's own lines, in the article's own notation, and no retained line is substituted or re-flowed.  Retained uncounted as the source-local context the proof needs: lines 250--256; lines 660--663; lines 676--678.  Nothing was omitted from any retained span, so the package carries no omission record.  No retained line contains a citation, so the wrapper carries no bibliography and no bibliography entry.  No retained line contains a figure environment, an image-inclusion command or drawing code, so no image is delivered and the package declares no asset dependency.  Editorial formatting only, in addition: an AMS \texttt{amsart} preamble that restates the article's own declarations and macros used by the retained lines, namely the environments \texttt{lemma}, \texttt{proposition}, \texttt{claim} on separate counters, together with the standard packages those lines need.  The retained environments print in this document as Lemma 1, Proposition 1, Claim 1 in order of appearance, whereas the article prints the same items with the numbering of its own full document.  The complete native source of the article as distributed in the arXiv e-print for this exact version is bundled unchanged as \texttt{sources/203292/source0.tex}.
\begin{lemma}\label{lem:properties of deg(g)}
    Let $p \in \widetilde{C}$, and let $u,v \in \L_C$ be linearly independent elements. Then the following hold.
    \begin{enumerate}
        \item If $\deg_p(u) \neq \deg_p(v)$, then $\deg_p(\alpha u + \beta v) = \max(\deg_p(u),\deg_p(v))$ for all $\alpha, \beta \in \kk^*$.\label{item:maximum condition}
        \item If $\deg_p(u) = \deg_p(v)$, then there exist $\alpha,\beta \in \kk^*$ such that $\deg_p(\alpha u + \beta v) < \deg_p(u)$.\label{item:decrease equal degree}
    \end{enumerate}
\end{lemma}

\begin{lemma}\label{lem:poles can be arbitrarily large}
    Let $p \in \widetilde{D} \setminus D$. Then there exists $N \in \ZZ$ such that
    $$\ZZ_{\geq N} \subseteq \deg_p(\g).$$
\end{lemma}

\begin{lemma}\label{lem:degrees at different points are different}
    Let $p,q \in \widetilde{D} \setminus D$ with $p \neq q$. Then there exists $w \in \g$ such that $\deg_p(w) \neq \deg_q(w)$.
\end{lemma}

\begin{proposition}\label{prop:dichotomy of poles at p and q}
    Let $p,q \in \widetilde{D} \setminus D$ with $p \neq q$. Then at least one of the following holds:
    \begin{enumerate}
        \item There exists $w \in \g_p'$ with a pole at $q$.
        \item There exists $w \in \g_q'$ with a pole at $p$.
    \end{enumerate}
\end{proposition}

\begin{proof}
    Assume, for a contradiction, that no elements of $\g_p'$ and $\g_q'$ have poles at $p$ or $q$. Given $n \in \deg_p(\g)$, define
    $$S_n \coloneqq \{\deg_q(w) \mid w \in \g \text{ and } \deg_p(w) = n\}.$$
    Suppose $S_n$ is unbounded below for some $n \in \deg_p(\g)$. Let $v \in \g$ such that $\deg_p(v) \geq 2 - n$ (which exists by Lemma~\ref{lem:poles can be arbitrarily large}), and let $u \in \g$ such that $\deg_p(u) = n$ and $\deg_q(u) \leq -2 - \deg_q(v)$ (which exists by the assumption that $S_n$ is unbounded below). Then
    $$\deg_p([u,v]) = \deg_p(u) + \deg_p(v) \geq 2, \qquad \deg_q([u,v]) = \deg_q(u) + \deg_q(v) \leq -2,$$
    so that $w \coloneqq [u,v] \in \g_q'$ with $\deg_p(w) \geq 2$ (so $w$ has a pole at $p$), a contradiction. It follows that $S_n$ is bounded below for all $n \in \deg_p(\g)$.

    Define a function
    \begin{align*}
        f \colon \deg_p(\g) &\to \deg_q(\g) \\
        n &\mapsto \min(S_n).
    \end{align*}
    Given $n \in \deg_p(\g)$, we let $w_n \in \g$ such that $\deg_p(w_n) = n$ and $\deg_q(w_n) = \min(S_n)$. We will prove that $f$ is the identity map through a series of claims, at which point we will be able to use Lemma~\ref{lem:degrees at different points are different} to prove the result.

    \begin{claim}\label{claim:f is injective}
        $f$ is injective.
    \end{claim}
    \begin{proof}[Proof of Claim \ref{claim:f is injective}]
        Assume, for a contradiction, that $f(n) = f(m)$ for some $n > m$. Then $\deg_q(w_n) = \deg_q(w_m)$, so Lemma~\ref{lem:properties of deg(g)} implies that there exists $v \in \spn\{w_n,w_m\}$ with
        $$\deg_p(v) = n, \qquad \deg_q(v) < \deg_q(w_n) = f(n).$$
        This contradicts the definition of $f$.
    \end{proof}

    \begin{claim}\label{claim:f is sub-additive}
        $f(n + m) \leq f(n) + f(m)$ for all $n,m \in \deg_p(\g)$ with $n \neq m$.
    \end{claim}
    \begin{proof}[Proof of Claim \ref{claim:f is sub-additive}]
        For $n,m \in \deg_p(\g)$ with $n \neq m$, we have
        \begin{align*}
            \deg_p([w_n,w_m]) &= \deg_p(w_n) + \deg_p(w_m) = n + m, \\
            \deg_q([w_n,w_m]) &= \deg_q(w_n) + \deg_q(w_m) = f(n) + f(m),
        \end{align*}
        where we used that $\deg_q(w_n) \neq \deg_q(w_m)$ by Claim \ref{claim:f is injective}. It follows that $f(n + m) \leq f(n) + f(m)$ by definition of $f$.
    \end{proof}

    Using Lemma~\ref{lem:poles can be arbitrarily large}, choose $N \geq 2$ such that $\ZZ_{\geq N} \subseteq \deg_p(\g)$.

    \begin{claim}\label{claim:f is increasing}
        $f(n) \geq n$ for all $n \in \ZZ_{\geq N}$.
    \end{claim}
    \begin{proof}[Proof of Claim \ref{claim:f is increasing}]
        Assume, for a contradiction, that $f(n) < n$ for some $n \in \ZZ_{\geq N}$. For $t \in \ZZ_{\geq 2}$, define
        $$X_t \coloneqq \{b n + r \mid 2 \leq b \leq t, 1 \leq r \leq \min(n,b - 1)\}.$$
        For $2 \leq b \leq n + 1$, there are $b - 1$ integers $m$ that can be written as $m = b n + r$ for some $1 \leq r < b$. Therefore,
        $$|X_{n + 1}| = 1 + 2 + \dots + n = \frac{n(n + 1)}{2}.$$
        For $b \geq n + 2$, there are $n$ integers $m$ that can be written as $m = b n + r$ for some $1 \leq r \leq n$, and thus
        $$|X_t| = \frac{n(n + 1)}{2} + (t - (n + 1))n = n t - \frac{n(n + 1)}{2}$$
        for all $t \in \ZZ_{\geq n + 1}$.
        
        Let $m = b n + r \in X_t$ for some $t \in \ZZ_{\geq 2}$, where $2 \leq b \leq t$ and $1 \leq r \leq \min(n, b - 1)$. Notice that $m$ can be written as
        $$m = k n + r(n + 1),$$
        where $k \coloneqq b - r > 0$. Applying Claim \ref{claim:f is sub-additive} inductively, we have
        $$f(m) = f(b n + r) = f(k n + r (n + 1)) \leq k f(n) + r f(n + 1).$$
        Letting $c \coloneqq f(n + 1)$ and $d \coloneqq f(n) - f(n + 1)$ (so that $f(n) = c + d$), we deduce that
        $$f(m) \leq k(c + d) + rc = (k + r)c + kd = b c + kd.$$
        Define
        $$a \coloneqq \begin{cases}
            d, &\text{if } d > 0, \\
            nd, &\text{if } d < 0.
        \end{cases}$$
        Since $1 \leq r \leq n$ and $k = b - r$ by definition, we see that
        $$b - n \leq k \leq b - 1.$$
        If $d > 0$, then $k \leq b - 1$ implies that $kd \leq (b - 1)d$, so we deduce from the above that
        $$f(m) \leq bc + kd \leq bc + (b - 1)d = (c + d)b - d = f(n)b - a \leq (n - 1)b - a,$$
        since $f(n) \leq n - 1$ by assumption. If $d < 0$, then $k \geq b - n$ implies that $kd \leq (b - n)d$, so we similarly deduce that
        $$f(m) \leq bc + kd \leq bc + (b - n)d = (c + d)b - nd = f(n)b - a \leq (n - 1)b - a,$$
        where we once again used that $f(n) \leq n - 1$. In either case, we have $f(m) \leq (n - 1)b - a$. Therefore, we see that
        $$f(m) \leq (n - 1)t - a$$
        for all $t \in \ZZ_{\geq 2}$ and $m \in X_t$.

        Notice that $\deg_q(w_n) \geq -1$ for all $n \in \deg_p(\g) \cap \ZZ_{\geq 2}$, since no element of $\g_q'$ can have a pole at $p$ by assumption. In particular, this means that $f(\ZZ_{\geq N}) \subseteq \ZZ_{\geq -1}$. Define
        $$Y_t \coloneqq \{k \in \ZZ \mid -1 \leq k \leq (n - 1)t - a\}.$$
        Let $t$ be a large enough integer such that
        $$nt - \frac{n(n + 1)}{2} > (n - 1)t - a + 2.$$
        Then the restriction
        $$\restr{f}{X_t} \colon X_t \to Y_t$$
        is an injective function. This is a contradiction, because
        $$|X_t| = nt - \frac{n(n + 1)}{2} > (n - 1)t - a + 2 = |Y_t|.$$
        It follows that $f(n) \geq n$ for all $n \in \ZZ_{\geq N}$.
    \end{proof}

    \begin{claim}\label{claim:f is the identity}
        $f(n) = n$ for all $n \in \deg_p(\g)$.
    \end{claim}
    \begin{proof}[Proof of Claim \ref{claim:f is the identity}]
        Let $g \colon \deg_q(\g) \to \deg_p(\g)$ be the function one gets by swapping the roles of $p$ and $q$ above. Then Claims \ref{claim:f is injective}--\ref{claim:f is increasing} are true if we replace $f$ by $g$. More precisely, use Lemma~\ref{lem:poles can be arbitrarily large} to get an integer $N' \geq 2$ such that $\ZZ_{\geq N'} \subseteq \deg_q(\g)$. Then $g(n) \geq n$ for all $n \in \ZZ_{\geq N'}$, by Claim \ref{claim:f is increasing} applied to $g$.

        Recalling that $\deg_q(w_n) = f(n)$ and $\deg_p(w_n) = n$, it follows that
        \begin{equation}\label{eq:g(f(n)) is less than n}
            g(f(n)) = \min\{\deg_p(w) \mid w \in \g \text{ and } \deg_q(w) = f(n)\} \leq n
        \end{equation}
        for all $n \in \deg_p(\g)$. Let $M \coloneqq \max(N,N')$. Given $n \in \ZZ_{\geq M}$, we have $f(n) \geq n \geq M$, and thus
        $$n \geq g(f(n)) \geq f(n) \geq n$$
        for all $n \in \ZZ_{\geq M}$, where the first inequality follows from \eqref{eq:g(f(n)) is less than n}, the second and third ones follow from Claim \ref{claim:f is increasing} applied to $g$ and $f$, respectively. It follows that $f(n) = n$ for all $n \in \ZZ_{\geq M}$.

        Now $n \in \deg_p(\g)$, and let $m \in \deg_p(\g)$ such that $m \geq M$ and $n + m \geq M$. Applying the above, we get $f(n + m) = n + m$ and $f(m) = m$, and thus
        $$n + m = f(n + m) \leq f(n) + f(m) = f(n) + m,$$
        where we used Claim \ref{claim:f is sub-additive} to deduce that $f(n + m) \leq f(n) + f(m)$. It follows that $f(n) \geq n$ for all $n \in \deg_p(\g)$. Similarly, we also get that $g(n) \geq n$ for all $n \in \deg_q(\g)$. We conclude that
        $$n \geq g(f(n)) \geq f(n) \geq n$$
        for all $n \in \deg_p(\g)$, from which it follows that $f(n) = n$ for all $n \in \deg_p(\g)$.
    \end{proof}

    By the definition of $f$, Claim \ref{claim:f is the identity} implies that $\deg_p(w) \leq \deg_q(w)$ for all $w \in \g$. Swapping the roles of $p$ and $q$, we also get that $\deg_p(w) \geq \deg_q(w)$ for all $w \in \g$. In other words, we have $\deg_p(w) = \deg_q(w)$ for all $w \in \g$. This contradicts Lemma~\ref{lem:degrees at different points are different}, which concludes the proof.
\end{proof}

\end{document}
